% Source: 13 - Mon Oct 26 2026.pdf, page 10. Content remains in source order.
\sourcepage{13}{10}
\section{Approximation analysis of greedy set cover}
\subsection{Two approximation bounds}
We will see two different arguments to upperbound $\rho$:
\begin{itemize}
\item Simple argument yielding $\rho(n)\leq\lceil\ln n\rceil$ with $n=|X|$.
\item Complex argument yielding $\rho(\langle X,\mathcal F\rangle)\leq\ln|S_{\max}|+1$ with
$|S_{\max}|=\max\{|S|:S\in\mathcal F\}$.
\end{itemize}
(NOTE: $|S_{\max}|$ can be $\ll n$!)

Interesting charging argument (uses weights).

\subsection{The logarithmic bound in the universe size}
\begin{notestheorem}
\textbf{THEOREM}\quad $\rho(n)\leq\lceil\ln n\rceil$, $n=|X|$.
\end{notestheorem}
PROOF. Call $U_t$ set $U$ at the start of iteration $t\geq1$. $U_t$ is the set of elements that still need to be covered.

Observe that the first time $\bar t$ when $|U_{\bar t}|=\varnothing$ implies that $|\mathcal C|=\bar t-1$.
% The inset repeats the algorithm; its repetition is retained.
\begin{notesalgorithm}
\textbf{APPROX-SET-COVER ($X,\mathcal F$)}
\begin{notespseudocode}
\pseudocodeline{0}{$U\gets X;\quad\mathcal C\gets\varnothing$}
\pseudocodeline{0}{\textbf{while} $(U\neq\varnothing)$ \textbf{do}}
\pseudocodeline{1}{$S\gets\operatorname{argmax}\{|T\cap U|:T\in\mathcal F\}$}
\pseudocodeline{1}{$U\gets U-S;\quad\mathcal C\gets\mathcal C\cup\{S\}$}
\pseudocodeline{0}{\textbf{return} $\mathcal C$}
\end{notespseudocode}
\end{notesalgorithm}
